\FloatBarrier
\section{FT flange}
\label{sec:ft-flange}

The twelfth beamline mesh study renames \texttt{flange-assembly} to \texttt{FTFlange} in four configurations:
runs 5874, 11620, 18304, and 20000. The stainless-steel material, passive 180-degree y rotation, and 1273.27
mm z translation remain unchanged. Both repositories and the GEMC2 CAD database use the new name. GEMC3 SQLite
banks remain in \texttt{default}, selected by run number. The original is archived, and all three installed
meshes are byte-identical.

The original is already a clean closed annulus with only 552 triangles, one component, and genus one.
Topology-preserving quadric decimation reduces it to 540 triangles (2.1739\% fewer); stored bytes fall from
27,684 to 27,084 (2.1673\% fewer). No holes are filled or surfaces reconstructed. The central passage and
overall bounds are retained. After edge collapse, vertices are restored to the two exact engineering face
planes, local z = -526.342956543 and -519.992980957 mm. The resulting STL has no boundary edges, non-manifold
elements, detected self-intersections, or zero-area triangles, and the builder reproduces it byte for byte.

Steel volume changes from 229348.479344 to 229348.228841 mm$^3$ (-0.000109224\%). Bidirectional vertex, edge,
and surface sampling gives a maximum deviation of 0.049822 mm, below the 0.05 mm guard. This sampled maximum
is not a proven global error bound. A 536-triangle trial exceeds that guard; larger reductions are rejected.
The small reduction reflects the already-coarse engineering tessellation. Table~\ref{tab:ft-flange-mesh}
summarizes the mesh changes, and Figure~\ref{fig:ft-flange} compares the stored triangles with a matched
camera.

The isolated benchmark selects run 5874 with \texttt{default}, using matched 5 GeV electron, proton,
positive-pion, and neutron primaries from the target, uniform in solid angle over 1.4--3.8 degrees and full
azimuth. Physics is FTFP\_BERT with no field or optical processes, one thread, default production cuts, and no
GUI or event output. Five alternating-order repetitions use 2,000 events per particle and seeds 12345--12349.
The executable matches the recent beamline studies, and all primary hashes agree within seed pairs.
Table~\ref{tab:ft-flange-tracking} reports median Geant4 beamOn wall times excluding initialization, with
original/optimized ratios and the ranges of paired ratios.

Tracking is effectively unchanged: median ratios are 1.0018 for electrons, 1.0001 for protons, 1.0003 for
positive pions, and 0.9993 for neutrons. Every paired range includes one; only three or four of the five pairs
are faster. These variations do not establish a tracking-speed gain. All 52 isolated processes complete
without Geant4 exceptions. Median loading process wall time changes from 0.784696 to 0.791953 s, CPU time
remains 0.72 s, and engine initialization wall time changes from 0.681824 to 0.672900 s. Startup peak RSS is
115.391 versus 115.375 MiB, effectively unchanged. Complete process CPU, wall time, memory, and per-seed
results are retained separately from the transport measurements.

The neighboring inner and outer assemblies already contact or overlap the flange before this change.
Independent triangle collision counts are 449 versus 449 for the inner assembly and 600 versus 583 for the
outer assembly; those tessellation-dependent counts are not overlap volumes or clearance proofs. The unchanged
inner assembly is not closed. A separate three-part 10,000-point Geant4 check reports one GeomSolids1001 and
eight GeomVol1002 exception blocks both before and after optimization. Sampled penetration depths reach about
0.249 mm. This mesh simplification retains the existing engineering interfaces; the assembly defects remain.

All four complete affected configurations construct 39 Geant4 volumes including the world. Both versions
produce two GeomSolids1001 and 21 GeomVol1002 blocks per complete construction. All 17 beamline configurations
pass volume-count construction checks, and the comparator passes 15 native geometry, 15 material, and 14 CAD
comparisons. Direct database checks allow only the four names and STL paths to change; every other mesh
retains its previous hash. The clean isolated benchmark does not certify complete geometry or
detector-response equivalence.

\begin{table}[htbp]
\centering
\caption{FT flange before and after conservative decimation.}
\label{tab:ft-flange-mesh}
\setlength{\tabcolsep}{6pt}
\begin{tabular}{lrr}
\toprule
Quantity & Original & Optimized \\
\midrule
Triangles & 552 & 540 \\
Bytes & 27,684 & 27,084 \\
Boundary edges & 0 & 0 \\
Components / genus & 1 / 1 & 1 / 1 \\
Volume (mm$^3$) & 229348.479344 & 229348.228841 \\
\bottomrule
\end{tabular}
\end{table}

\begin{table}[htbp]
\centering
\caption{Isolated FT flange tracking; five repetitions of 2,000 events per particle.}
\label{tab:ft-flange-tracking}
\setlength{\tabcolsep}{6pt}
\begin{tabular}{lrrrr}
\toprule
Particle & Original (s) & Optimized (s) & Ratio & Paired ratio range \\
\midrule
Electron & 2.68847 & 2.68375 & 1.0018 & 0.9998--1.0054 \\
Proton & 3.61487 & 3.61435 & 1.0001 & 0.9957--1.0029 \\
Positive pion & 3.06986 & 3.06909 & 1.0003 & 0.9994--1.0033 \\
Neutron & 1.75223 & 1.75338 & 0.9993 & 0.9986--1.0051 \\
\bottomrule
\end{tabular}
\end{table}

\begin{figure}[p]
\centering
\includegraphics[width=\linewidth]{figures/ft_flange_stl_comparison.pdf}
\caption{Original and optimized FT flange meshes with identical bounds and camera. The central passage, flat
axial faces, and overall shape are retained; the already-coarse tessellation changes modestly.}
\label{fig:ft-flange}
\end{figure}
